% ****** Start of file apssamp.tex ******
%
%   This file is part of the APS files in the REVTeX 4.2 distribution.
%   Version 4.2a of REVTeX, December 2014
%
%   Copyright (c) 2014 The American Physical Society.
%
%   See the REVTeX 4 README file for restrictions and more information.
%
% TeX'ing this file requires that you have AMS-LaTeX 2.0 installed
% as well as the rest of the prerequisites for REVTeX 4.2
%
% See the REVTeX 4 README file
% It also requires running BibTeX. The commands are as follows:
%
%  1)  latex apssamp.tex
%  2)  bibtex apssamp
%  3)  latex apssamp.tex
%  4)  latex apssamp.tex
%
\documentclass[%
reprint,
superscriptaddress,
%groupedaddress,
%unsortedaddress,
%runinaddress,
%frontmatterverbose, 
%preprint,
%preprintnumbers,
%nofootinbib,
%nobibnotes,
%bibnotes,
 amsmath,amssymb,
 aps,
 prx,
%pra,
%prb,
%rmp,
%prstab,
%prstper,
floatfix,
]{revtex4-2}

\usepackage[english]{babel}

\usepackage[hidelinks]{hyperref}% add hypertext capabilities
\usepackage{graphicx}% Include figure files
\usepackage{dcolumn}% Align table columns on decimal point
\usepackage{bm}% bold math

\usepackage{comment}
\usepackage{physics}
\usepackage{cleveref}
\usepackage{xfrac}
\usepackage{xcolor}
\usepackage{booktabs}
\usepackage{tabularx}
\usepackage{threeparttable}


%% CUSTOM MACROS
\newcommand{\unsure}[1]{{\textcolor{red}{#1}}}
\newcommand{\ketlpOI}{\ensuremath{\ket{\text{LP}_{01}}}}
\newcommand{\ketlpIIa}{\ensuremath{\ket{\text{LP}_{11\text{a}}}}}
\newcommand{\ketlpIIb}{\ensuremath{\ket{\text{LP}_{11\text{b}}}}}

\newcommand{\prob}[1]{\ensuremath{\operatorname{Pr}\left(#1\right)}}
\newcommand{\probgiven}[2]{\ensuremath{\operatorname{Pr}\left(#1 \middle| #2\right)}}
\newcommand{\probguess}[1]{\ensuremath{\operatorname{Pr}_g\left(#1\right)}}

\newcommand{\linbo}{{$ \text{LiNbO}_3 $} }
\newcommand{\guilherme}[1]{\textcolor{blue}{[guilherme: #1]}}
\usepackage{float}

\begin{document}

\preprint{APS/123-QED}

\title{An optical-fibre-integrated buffer for packet-switched quantum networks} 


% \author{Author name(s)}
% \address{Author affiliation and full address}
% \email{e-mail address}
%%Uncomment the following line to override copyright year from the default current year.
%\copyrightyear{2024}

\author{Daniel~Spegel-Lexne}
\email{daniel.spegel-lexne@liu.se}
\affiliation{Institutionen för Systemteknik, Linköpings Universitet, 581 83, Linköping, Sweden}

\author{Joakim~Argillander}
\affiliation{Institutionen för Systemteknik, Linköpings Universitet, 581 83, Linköping, Sweden}

\author{Martin~Clason}
\affiliation{Institutionen för Systemteknik, Linköpings Universitet, 581 83, Linköping, Sweden}

\author{Åsa~Claesson}
\affiliation{Fibrelab, RISE-Research Institutes of Sweden, Fibrevägen 2-6, 82450, Hudiksvall, Sweden}

\author{Kenny~Hey~Tow}
\affiliation{Fibre Optics, RISE-Research Institutes of Sweden, Electrum 236, 16440, Kista, Sweden}

\author{Gustavo~Lima}
\affiliation{Departamento de F\'{\i}sica, Universidad de Concepci\'on,
             160-C Concepci\'on, Chile}
\affiliation{Millennium Institute for Research in Optics,
             Universidad de Concepci\'on, 160-C Concepci\'on, Chile}

\author{João~M.~B.~Pereira}
\affiliation{Fibre Optics, RISE-Research Institutes of Sweden, Electrum 236, 16440, Kista, Sweden}

\author{Guilherme~B.~Xavier}
\email{guilherme.b.xavier@liu.se}
\affiliation{Institutionen för Systemteknik, Linköpings Universitet, 581 83, Linköping, Sweden}

\date{\today}% It is always \today, today,
             %  but any date may be explicitly specified

\begin{abstract}

Packet-switched quantum networks require buffers that can delay qubit payloads while routing information is read out in real time. Previous approaches have not provided this functionality in a fully fibre-integrated architecture compatible with telecom infrastructure. Here we demonstrate an optical-fibre-integrated buffer, based on a recirculating loop and a fibre storage line, in which the storage time of a polarisation-encoded qubit payload is determined by readout of an attached packet header. The key component behind this achievement is an ultra-low-loss poled fibre phase modulator, which provides fast, polarisation-insensitive switching directly in fibre and allows header and payload to be processed within the buffer. We demonstrate storage and retrieval of polarisation-encoded qubit payloads for storage times up to 47 $\mu$s, with an average quantum bit error rate of 1.8\% together with stable operation over several hours. These results establish a practical fibre-based architecture for packet-level quantum network buffering that can easily integrate into the current telecommunication infrastructure opening up new paths for deployment of the quantum internet.



\end{abstract}

\maketitle

\section*{Introduction}

Realizing large-scale quantum networks, often framed as a “quantum internet”, would enable functionalities that are fundamentally out of reach for classical communication networks, including device-independent security and distributed quantum computing \cite{kimble2008quantum,wehner2018quantum}. Essentially, in all architectures, photons in the telecom band are the natural carriers of quantum information across metropolitan and intercity links owing to the low-loss properties of commercial optical fibres in the 1550 nm spectral region \cite{gisin2007quantum, pirandola2020advances}. Quantum key distribution (QKD) is a well-established quantum communication application, having had many demonstrations ranging from foundational experiments to modern, ultra-long-distance and side-channel attack secure versions \cite{bennett1984bb84,ekert1991bell,lo2012mdi,TF2018}, as well as field demonstrations now spanning terrestrial fibre networks and free-space and satellite links \cite{liao2017satelliteqkd,yin2017satelliteentanglement, zhou2024combfield, pittaluga2025deployed}. As these systems evolve from point-to-point demonstrations to multi-user networks, they increasingly require network-layer primitives, such as packet buffering, without degrading the fragile quantum states \cite{wehner2018quantum,pompili2022networkstack}.

Recently, significant effort has been dedicated to developing a framework for packet-switched quantum networks as a natural extension from classical packet-switched communication extensively used over the Internet \cite{DiAdamo2022, Mandil2023, Yoo2024}.  Within a network node, quantum memories \cite{lvovsky2009memory,heshami2016memories,bussieres2013applications} are needed to act as buffers allowing time for header processing belonging to each quantum packet. Typical quantum memory implementations, although nowadays offering excellent performance in terms of storage lifetime or efficiency, have the general disadvantage of being relatively complex and not directly spectrally compatible with optical-fibre telecom infrastructure, as well as requiring cryogenic cooling and bulk optical elements making them challenging to integrate into real-world network nodes \cite{Ortu2022, Duranti2024}. 

Another approach to try to bridge the gap towards more practical quantum memories is based on the use of an optical loop (working as the quantum memory), combined with an optical switch for storage/retrieval \cite{Pittman2002}. This approach has been highly successful with many relevant applications being demonstrated \cite{Makino2016, Kaneda2017, Kaneda2019, Meyer-Scott2022, Hou2023, Weinbrenner2024, Hanamura2025, Pang2023, Bonsma-Fisher2023, Bonsma-Fisher:24}. All these previous approaches relied on bulk optical elements, making it difficult to integrate with fibre-optical components, thus greatly limiting their use for quantum communication and networking applications. Recently, efforts on fibre-based recirculating memories for quantum applications have included a fibre loop with an electro-optical switch with storage periods of several $\mu s$ \cite{FookLee2024}. Other very recent approaches managed to have very much faster storage periods but with the recirculating loop itself implemented in free-space \cite{Evans2023, Cheng2025, Sun2025}. Furthermore, none of these previous experiments have demonstrated higher-level functions like processing a packet with a payload containing quantum states. Therefore, an efficient, fully fibre integrated approach with fast temporal responses capable of packet handling is still missing. 

Here we cover this gap by demonstrating an optical-fibre-integrated processing buffer for packet-switched quantum networks, based on a Sagnac-loop switch \cite{alarcon2020sagnacswitch} and a fibre storage line controlled by an ultra-low-loss poled-fibre phase modulator \cite{myers1991poledsilica,kazansky1997poling,fujiwara1995uvpoling,fokine2002mzswitch, Spegel-Lexne2025}. This key device exhibits much lower losses compared with commercial fibre-optical lithium niobate (LiNbO$_3$) modulators, as well as exhibiting polarisation independence, a key property for handling polarisation-encoded quantum states.  We successfully demonstrate storage, and subsequent retrieval of a polarisation qubit payload across the fibre buffer, as well as active readout of a header while the payload remains stored. Our results demonstrate a practical and fully telecom compatible fibre architecture for processing buffers for the quantum internet. Parallel to this work in a companion submission \cite{Guerrero2026}, a completely distinct application of poled fiber modulators for quantum information was demonstrated, through the construction of an ultra-high efficiency active receiver for high-dimensional photonic quantum states in a quantum key distribution demonstration. Together we demonstrate the first uses of poled fibre phase modulators for quantum information processing applications, whose unique characteristics, namely, ultra-low-loss, fast response times, polarisation state insensitivity and direct fibre compatibility may open the path forward for highly demanding quantum communication applications. 

\section*{Results}

\subsection*{Quantum network buffering}

The quantum internet requires switching and real-time processing of packets of quantum states at the network nodes \cite{DiAdamo2022, Mandil2023, Yoo2024}, in a similar fashion to the standard internet. A typical packet structure consists of a payload containing a number of quantum states and a classical header, which holds classical information needed for real-time processing of the packet at each network node (Fig. \ref{fig:buffering}a). We encode the header and payload at the same wavelength, with our implementation being able to split the header from the payload within the buffer, owing to the high-speed response of the poled phase modulator.

A typical quantum networking scenario occurs in a star network, where multiple users are connected to a central node, responsible for processing and rerouting the packets (Fig. \ref{fig:buffering}b). As an example, $N$ simultaneous users wish to distribute quantum states to $M$ other users, all connected through the central processing node, with a multi-port switch that can connect any input to any output (Fig. \ref{fig:buffering}c). At the end of each quantum channel between the $N$ transmitting users and the central node, packet buffers are deployed to appropriately delay the packets before they are processed by the switch. The quantum states are sequentially prepared within a payload by each user and a header containing destination information (i.e. an address) is added to each payload. The header will be processed at the central node, to determine how to configure the switch for the correct routing operation. The node further checks whether the continuing communication channels to each of the $M$ destination users is available. If packets with the same final destination (i.e. $M_1$) arrive at the node with a time difference $\Delta\tau < \tau_p$, where $\tau_p$ is the time duration of the packet, then they will overlap in time in the outgoing channel, and the receiver will be unable to distinguish between them. Buffers are then used before the switch to delay the packets appropriately, such that they can all enter the same channel to $M_1$ without overlapping in time. 

\begin{figure}[ht!]
    \centering
    \includegraphics[width=0.5\textwidth]{figures_fontsync/Fig1.pdf}
    \caption{Quantum network packet switching. a) Packet structure consisting of a header containing classical information and a payload consisting of quantum states. b) A star network configuration, with multiple users connected through a central switching node. c) $N$ users are connected to $M$ users in a star configuration, where the central node has buffering and switching functionality. Packets can be sent from any of the $N$ to the $M$ users at will. As an example, if packets from users $N_1$, $N_2$ and $N_3$, which are all addressed to the same end user ($M_1$), arrive simultaneously at the central node, then the packets will collide in the outgoing communication channel to $M_1$. Buffering is therefore needed to time-multiplex the packets to avoid collisions in the outgoing link. The central node also needs to process the header, which contains source and destination information for instance. With this information, the delay amount in the buffer is decided by the node. Our implementation allows both header processing and buffering functionality for a packet. 
   }
    \label{fig:buffering}
\end{figure}

\subsection*{Poled optical fibre phase modulator}

Previous active loop recirculating memories often rely on bulk electro-optical modulators for storage/retrieval control, which limits integration capabilities within optical fibre communication systems. Commercial fibre pig-tailed electro-optical modulators have typically high insertion losses ($\sim 50\%$) that greatly hamper their use for multi-pass quantum photonics applications or measurements that have a high minimum threshold detection efficiency. Therefore we resort to optical fibre poling technology to construct an in-fibre low-loss phase modulator for the recirculating loop memory. The poled fibre phase modulator used here was fabricated from a specially designed single-mode silica fibre containing two longitudinal holes running parallel to the core. These holes were drilled at the preform stage before fibre drawing and subsequently filled with bismuth electrodes in a pressure chamber, producing the cross-section shown in Fig. \ref{poledmodulator}a \cite{fokine2002mzswitch}.


\begin{figure}[ht!]
    \centering
    \includegraphics[width=0.45\textwidth]{figures_fontsync/Fig2.pdf}
    \caption{Poled fibre modulator. a) The cross-section of the poled fibre phase modulator is shown, imaged using a Vytran glass processor. b) Interference curves at the two outputs of a fibre-optical Sagnac interferometer driven by the poled fibre phase modulator. The outputs are measured with the two single-photon detectors. The interferometer is the same one employed in our experiment, please see in the next section for details. The error bars represent the standard deviation assuming a Poissonian photon number distribution. c) Time domain response from a driving pulse applied to the poled modulator within the Sagnac interferometer. In this measurement, the classical optical output is measured with a p-i-n photo-receiver. The rise and fall times are measured to be 22.4 ns and 26.4 ns respectively.
   }
    \label{poledmodulator}
\end{figure}

 The application of an electric field between the electrodes causes a phase shift to be applied to the propagating light due to the Kerr effect. Due to the low optical non-linearity of glass, the optical fibre is first poled by coupling to its core 532 nm light from a Nd:YAG laser, while simultaneously applying a high-voltage of $\approx5$ kV, allowing electric charges to redistribute and leaving a long-lived recorded electric field in the glass after the poling field is removed \cite{Camara15}. After poling, the needed voltage to apply a $\pi$ phase shift decreases from approximately ($V_{\pi}\approx5$ kV) for 950 V to our device (Fig. \ref{poledmodulator}b), and in terms of time response, we obtain a rise and fall time of 22.4 and 26.4 ns respectively (Fig. \ref{poledmodulator}c). Since the phase modulator is essentially just an optical fibre, the losses are very low and mainly limited by the fibre connectors, which we measure to be 0.4 dB.

\subsection*{Fibre-optical active loop memory}

A recirculating quantum memory typically consists of a resonating circuit, with an active device needed to inject the photonic quantum states onto the circuit, and then retrieve it after some chosen time \cite{Pittman2002}. We employ a combination of a fibre-optical Sagnac interferometer with a fibre delay line and a mirror to complete the resonator circuit. The Sagnac interferometer acts as an optical switch, determining whether the states are stored or retrieved. Such interferometers are very fast optical switches with several applications in quantum information \cite{Alarcon2023, Spegel-Lexne2024, Raj2025, argillander2025, Vijayadharan2025}. Our experimental setup is depicted in Fig. \ref{fig:placeholder}.



\begin{figure*}[ht!]
    \centering
    \includegraphics[width=1.0\textwidth]{figures_fontsync/Fig3.pdf}
    \caption{All-fibre quantum packet buffer implementation. We prepare weak coherent polarisation encoded qubits in the state preparation stage. A fibre coupler is employed at the output of this stage, and connected to a superconducting nanowire single-photon detector (SNSPD) for photon number calibration purposes. Then an optical circulator leads to the routing stage consisting of a Sagnac interferometer with the poled fibre phase modulator used to control whether a relative $\phi=0$ or $\phi=\pi$ phase shift is applied onto the packets, thus causing it to be routed to the storage line or to exit the loop back through the circulator. A fibre delay is employed in the Sagnac loop to ensure the phase shift is applied to the entire packet in only one of the counter-propagating directions. The storage line consists of a 100 m fibre delay, a manual polarisation controller and a mirror composed of a fibre Bragg grating. The measurement stage consists of a manual polarisation controller and a polarising beam splitter, whose two outputs are connected to two other SNSPDs. Att: Optical attenuator;  Circ: Optical circulator; FBG: Fibre Bragg grating; HWP: Half-wave plate; IM: Intensity modulator; PC: Manual polarisation controller; PFPM: Poled fibre phase modulator; Pol: Linear polariser; VOPC: Voltage operated polarisation controller; PBS: Polarising beam splitter; SNSPD: Superconducting nanowire single-photon detector.}
    \label{fig:placeholder}
\end{figure*}

A source of weak coherent states, comprising of a continuous wave semiconductor laser diode at 1546.9 nm in tandem with an LiNbO$_3$ intensity modulator, produces a payload of 16 sequential pulses, with a period of 20 ns and a width of 10 ns at a repetition frequency of 975 Hz. A high voltage DC source feeds a high-voltage switch (Berkeley Nucleonics) capable of delivering up to 1 kV of pulse amplitude. As the modulator acts as a capacitor, it has very high impedance. Thus a 50 $\Omega$ high-voltage power resistor is connected in parallel to the modulator. This limits the repetition rate / pulse width product as the drawn current is limited by the DC source. The polarisation states of the pulses are set with a linear polariser followed by a half-wave plate (HWP). The pulses are then attenuated to the single photon level, (average 0.1 photons per pulse), approximately $1.6$ photons per payload.

The payload is then delivered to the fibre-optical Sagnac interferometer, which consists of a circulator, a 50 : 50 fibre coupler, a 1 km spooled fibre delay, the poled-fibre phase modulator and a manual polarisation controller to maximise the interference of the superposing co- and counter-propagating paths back on the fibre coupler. In one of the Sagnac outputs, the storage line is connected, consisting of a 100 m delay fibre, a manual polarisation controller and a fibre Bragg grating (FBG) acting as a mirror. The lengths of the fibre spools were chosen for convenience, and different lengths can be used depending on the desired payload length and overall storage time per cycle.

The other output, going back through the circulator, is then connected to a manual polarisation controller and a fibre polarisation beamsplitter (PBS) to perform a projective measurement on the retrieved polarisation qubits. Finally, the two outputs from the PBS are connected to two superconducting nanowire single-photon detectors (SNSPDs), within the same cryostat of the calibration detector. They have a nominal system detection efficiency $>$ 90\% and less than 30 dark counts per second for each detector. A time tagger (IdQuantique ID1000) is used to process the detection events. 

The Sagnac interferometer acts as a tunable optical switch, routing the photon between the two outputs depending on the setting $\phi$ of the poled fibre phase modulator. The relation between the applied phase difference and the probability of detecting a photon at the two outputs $D_1$ and $D_2$ of the interferometer is proportional to  $\textrm{cos}^2\left(\frac{\phi}{2}\right)$ and $\textrm{sin}^2\left(\frac{\phi}{2}\right)$ \cite{Alarcon2023}. If no relative phase is applied ($\phi=0$) the photon will exit through the same path it entered. If the phase is set to $\phi=\pi$ the photon will exit through the other output and thus, in our setup, enter the storage line. The photon is reflected back by the FBG and reenters the Sagnac loop. 

Since now the photon is entering from the opposing port, if $\phi=0$ the photon reflects back to the storage line and it becomes effectively trapped in a resonating circuit. Once a relative phase $\phi=\pi$ is applied, then the single-photon exits to the circulator and is retrieved from the memory. It then propagates to the measurement operation by the PBS before arriving at either detector $D_1$ or $D_2$.

The poled fibre phase modulator is driven by a high-voltage driver consisting of a DC power supply and a fast high-voltage switch. It generates two consecutive pulses separated in time, each 330 ns wide with 950 V amplitude, corresponding to a $\pi$ phase shift (inset Fig. \ref{poledmodulator}). The first pulse causes the qubit packet to be stored while the second retrieves it. By adjusting the delay between the two consecutive driving pulses, the photon payload can be stored and retrieved for a different number of cycles in the buffer. The period between each cycle where readout is possible is about $\approx 6 \mu s$ for the particular lengths of optical fibre in the Sagnac loop and storage line in our experiment. Note that the minimum storage time is one cycle.   


\subsection*{Experimental results}

We initially demonstrate the retrieval of a payload over different storage cycles, as determined by the delay between the high-voltage driving pulses. The payload consisting of the train of 16 weak coherent states is sent to the input of the Sagnac loop, with the measurement data acquired by accumulating the detection events with the time tagger with each cycle delay measured with a 30 sec integration time, and all the curves are superposed and plotted in Fig. \ref{fig:maxloopmeasurement}. The PBS before the detectors was removed for this initial measurement. The manual polarisation controllers are needed to compensate residual polarisation transformations for each storage cycle, and ensure optimal interferometric visibility at the Sagnac interferometer.

We clearly see the expected exponential decay in the detection probability between storage cycles, where the losses come from the fibre spools (0.3 dB), poled phase modulator (0.4 dB), fibre coupler (0.5 dB), FBG (0.1 dB), and additional losses from connectors and splices (1.2 dB), amounting to an efficiency of $\approx 56\%$ per cycle. There is also an additional loss of approximately 1 dB for each pass through the circulator, but this only happens once independently of the number of storage cycles. We can still observe detection events following retrieval after 10 storage cycles, which corresponds to approximately 60 $\mu s$. The first storage peak in the graph corresponds to no high-voltage driving pulse to the modulator ($\phi = 0$), and is used as a reference. We also plot in Fig. \ref{fig:maxloopmeasurement} the expected performance if a commercial LiNbO$_3$ phase modulator would be used instead (typical 3 dB loss).


\begin{figure}[ht!]
    \centering
    \includegraphics[width=\linewidth]{figures_fontsync/Fig4.pdf}
    \caption{Buffer efficiency. A payload of 16 weak coherent states (WCSs) is stored in the buffer, and subsequently retrieved. We employ 0.1 photons per state on average (1.6 photons per payload) on the input of the buffer, to demonstrate storage for typical weak coherent state amplitudes in practical applications. The different curves corresponding up to ten storage cycles are obtained individually for each delay of the retrieval high voltage pulse, and acquired with a time tagger. They are then all plotted superposed on the figure. The green data points represent the average detected single-photon counts for all the 16 WCSs for that particular storage loop, and the green curve is an exponential decay fit. We also show for reference in the red curve, the expected performance if a commercial LiNbO$_3$ phase modulator is used instead of the poled fibre modulator. The insets zoom in on the last three storage cycles showing the individual detected counts for each WCS.}
    \label{fig:maxloopmeasurement}
\end{figure}

We then experimentally demonstrate the storage and retrieval of polarisation-encoded weak coherent states. The PBS at the measurement stage is reinstalled, and we employ the half-wave plate $\theta$ in the beginning of the setup to generate the superposition $|\psi_{\textrm{input}}\rangle = \textrm{cos}(2\theta)|H\rangle + \textrm{sin}(2\theta)|V\rangle$, where $|H\rangle$ and $|V\rangle$ represent the horizontal and vertical polarisation states respectively. We then sweep $\theta$ and thus, for each input state, the counts at the PBS outputs are recorded over 30 s, for 5 different storage cycles, the last one corresponding to 47 $\mu$s. This experiment was done for the two mutually unbiased polarisation measurement bases used in QKD BB84-protocols \cite{bennett1984bb84}, $\mathbf{Z}$ and $\mathbf{X}$, which are changed using the manual polarisation controller before the PBS. The corresponding polarisation curves for the n$_{th}$ storage cycle are presented in figure \ref{fig:polstore}a. 

In Fig. \ref{fig:polstore}b, we show the probabilities to detect each of the four BB84 transmitted states by setting the input polarisation state angle to $\theta_{H,V} = \{0, 45\}$ and $\theta_{D,A} = \{22.5, 67.5\}$, where the subscripts $D$ and $A$ correspond to the superposition states $|D\rangle = 1 /\sqrt{2}(|H\rangle+|V\rangle)$ and $|A\rangle = 1 /\sqrt{2}(|H\rangle-|V\rangle)$. We then record the counts over the two measurement bases $\mathbf{Z}$ and $\mathbf{X}$, corresponding to the projections onto the $|H\rangle ; |V\rangle$ and $|D\rangle ; |A\rangle$ states respectively. The probability of detection for the $|i\rangle$ state is calculated as $P_{|i\rangle}=N_{D^{|i\rangle}}/(N_{D^{|i\rangle}}+N_{D^{|i\rangle\perp}})$, where $N_{D^{|i\rangle}}$ and $N_{D^{|i\rangle\perp}}$ correspond to the number of detections on the matching and on the orthogonal detector respectively, for the $|i\rangle$ input state. The quantum bit error rate (QBER) is then written as QBER = $1 - P_{|i\rangle}$, and we calculated it from the detection probabilities in Fig. \ref{fig:polstore}b for the different storage times respectively as $1.06 \pm 0.04\%$, $1.11 \pm 0.08\%$, $1.46 \pm 0.13\%$, $2.74 \pm 0.39\%$, $2.29 \pm 0.55\%$ respectively.



\begin{figure*}[ht!]
    \centering
    \includegraphics[width=\linewidth]{figures_fontsync/Fig5.pdf}
    \caption{Quantum information storage. a) Normalised single-photon detection following the projective polarisation measurement in the measurement stage for different storage times, up to a maximum $n=8$ corresponding to 47 $\mu s$. It is carried out as a function of the input half-wave plate angle ($\theta$) in the superposition $|\psi_{\textrm{input}}\rangle = \textrm{cos}(2\theta)|H\rangle + \textrm{sin}(2\theta)|V\rangle$. For each storage cycle we project the state $|\psi_{\textrm{input}}\rangle$ over the two mutually unbiased bases used in BB84 QKD, the computational ($\mathbf{Z}$) and the logical ($\mathbf{X}$) bases. The integration time for each data point is 30 secs. Error bars come from the assumption that each detection event is governed by Poissonian statistics. b) We show the corresponding detection probabilities for each storage cycle number, by preparing each of the four prepared BB84 states and projecting onto each of the same four states, simulating a QKD protocol for different storage times, with an average QBER of $1.85 \pm 0.13\%$ over all cycles.}
    \label{fig:polstore}
\end{figure*}

We next evaluate the long-term stability of the buffer by continuously monitoring the QBER for a fixed input state, ($|H\rangle$). The measurement sequence starts at the first storage cycle, then proceeds to cycles n=3 and n=6, before returning to n=1. For this stability measurement, the mean input photon number is increased to 1.0, since the focus is on temporal stability. The main source of instability is the residual time-dependent birefringence of the optical fibre, which slowly changes the polarisation state reaching the measurement PBS. For each storage cycle, data is acquired for 30 s. At the start of each measurement round, the computer connected to the time tagger calculates the current detection probability (P$_{|H\rangle}$). If P$_{|H\rangle}<95\%$, corresponding to a QBER above 5\%, the measurement is paused and an automatic recalibration routine is initiated. This recalibration is performed using an all-fibre voltage-operated polarisation controller (VOPC) with 0.2 dB insertion loss, which realigns the polarisation state of the retrieved weak coherent states with the PBS before data acquisition resumes. The procedure typically takes a few minutes, depending on the initial polarisation condition. The longest storage case is the most sensitive to residual birefringence, as expected, because of the longer fibre propagation distance. Such automatic polarisation stabilisation is routinely used in fibre-based experiments and commercial QKD systems. The results, shown in Fig.~\ref{fig:visstab}, demonstrate stable long-term operation over more than 12 hours.


\begin{figure}[ht!]
    \centering
    \includegraphics[width=0.5\textwidth]{figures_fontsync/Fig6.pdf}
    \caption{Long-term stability of the buffer. QBER for detecting the state $|H\rangle$ on basis $\mathbf{Z}$ as a function of time for different storage cycles. Error bars come from the assumption that each detection event is governed by Poissonian statistics.}
    \label{fig:visstab}
\end{figure}

Finally, to demonstrate the full functionality of the buffer, we show processing of a header determining how long the payload should be stored, emulating the scenario in Fig. \ref{fig:buffering}. The header is added as three extra pulses to the beginning of the 16 weak coherent state payload, allowing the encoding of up to eight different addresses, assuming direct digital amplitude modulation. The central node then decides on the delay to be applied for each packet, depending on the destination address. In our proof-of-principle implementation, we map each header value to a different delay in terms of the storage cycle number. The header has a much higher average photon number compared to the payload, ensuring the header is successfully read out by the SNSPDs. Note that in an intermediate node in a network, simpler gated-mode avalanche diode-based single-photon detectors could be used to read out the header, or even high-sensitivity photo-receivers if the header intensity is tuned to be sufficiently high. The full packet structure is created from an arbitrary waveform generator producing the desired driving pulses to the amplitude modulator at the output of the laser. 

The possibility of performing real-time packet processing requires fast response times given the cumulative decoherence or loss probability as a function of storage time in a quantum memory. The poled fibre phase modulator also satisfies this strict requirement, allowing us to split the header from the payload without any need for additional time separation between the header and payload. The results are shown in Fig. \ref{fig:header}, where we see the time-of-arrival events acquired from the time-tagger for each number of different storage cycles with the corresponding header already split. We show a copy of the full packet before reaching the buffer (red shaded region), using a 10 : 90 splitter (not shown for simplicity) and the SNSPD. The header is then split off from the payload (blue shaded region), which stays in the loop and proceeds for another round. The header information containing the 3-bit binary word corresponding to the desired delay sets the delay of the applied retrieval pulse. The different retrieval times for each payload are shown in the green area. 

\begin{figure*}[ht!]
    \centering
    \includegraphics[width=1.0\linewidth]{figures_fontsync/Fig7.pdf}
    \caption{Quantum packet processing. The packet consists of a header containing three strong pulses, and the payload consists of 16 weak coherent states. The panels from top to bottom correspond to storage cycles of $n = \{2, 3, 4, 6\}$ respectively. On the left the red shaded waveform shows a copy of the packet before entering to the memory. The insets zoom in on the full packet before processing (red region), then the header after being split off from the payload (blue region) and the payload after being appropriately delayed (green region). Please see the text for further details. The different encoding values for the header can be seen based on on/off encoding for each respective delay for the payload on the insets.}
   \label{fig:header}
\end{figure*}



Figure~\ref{fig6} places our buffer in the broader landscape of room-temperature quantum matter memories \cite{Reim2011, Finkelstein2018, Kaczmarek2018, Thomas2023, Thomas2024} and photonic recirculating loop platforms.  We compare the effective qubit yield, which we define as $Y = \eta\,N$, where $N$ is the number of stored qubits and $\eta$ the storage efficiency per loop (or equivalently the total efficiency in the case of matter memories \cite{Reim2011, Finkelstein2018, Kaczmarek2018, Thomas2023, Thomas2024}). It quantifies the expected number of qubits recovered per storage event. In our case, through the storage and retrieval of a $16$-qubit packet at $\sim\!56\%$ per-cycle efficiency, our buffer
reaches $Y\approx 9$ qubit yield, overcoming the best previous result by nearly an order of magnitude.



\begin{figure}[ht!]
    \centering
    \includegraphics[width=1.0\linewidth]{figures_fontsync/Fig8.pdf}
    \caption{Effective qubit yield of room-temperature quantum  memories and photonic buffers. For each platform we plot the effective qubit yield $Y=\eta N$, which represents the expected number of qubits retrieved per storage event vs. the maximum reported storage time.  Loop memories consist of active recirculating memories while matter memories are included as a comparison, since they allow continuous on-demand retrieval. For the loop memories, we compare the different efficiencies per loop pass, while for matter memories we report the total efficiency. The efficiencies per cycle for references \cite{Bonsma-Fisher2023, Bonsma-Fisher:24, FookLee2024} were not explicitly given and were thus inferred from the experimental data provided.}
    \label{fig6}
\end{figure}









\section*{Discussion}

As point-to-point quantum communication protocols get more sophisticated the focus begins to shift to networking applications. In spite of significant proof-of-principle experiments over the years to implement different networking protocols and primitives, a quantum buffering device capable of real-time packet processing is still missing. Furthermore, for such a buffer to be of more practical value, it should be easy to integrate with current telecom networking hardware.

Here we are able to demonstrate such a buffer. Our approach was done at the critical telecom spectral window, allowing full compatibility with current systems. We implement a loop memory based configuration using a Sagnac interferometer which controls the storage of the quantum states in a storage line comprised of a fibre delay and a fibre Bragg grating, operating as a mirror. A major differentiator of our implementation compared with previous demonstrations is the use of a poled fibre phase modulator, allowing our platform to be fully fibre-based with fast response times, critical for packet-processed quantum networks. Additionally, compared to typical commercial fibre pigtailed modulators, our device offers much lower insertion loss as well as low polarisation dependence, properties which have prevented previous experiments from reaching the stringent performance requirements for a practical quantum internet.  

These first results are very promising as they do not require cryogenic systems, or complex and sensitive interfaces between free-space and fibre-optics. Our proof-of-principle platform can be further improved in terms of losses, by approximately 1.1 dB per pass through the replacement of four FC/PC connectors inside the Sagnac loop with fibre splices (1.0 dB), and a fibre coupler with lower excess losses (0.1 dB), thus improving the overall transmissivity per storage cycle to 72$\%$. This value can be slightly different depending on the fibre delay values, which will be determined by the desired packet length. 

Since the buffer is designed to process entire packets, much denser payloads and headers are possible by using much narrower and closely packed pulse widths, without degrading the performance of the storage. For instance, recent results have shown the use of 45 ps wide pulses with 400 ps separation \cite{Grunenfelder2023}. The repetition rate of the buffer can be greatly improved by employing DC sources with higher current limits or a pulse generator capable of directly driving  high-impedance capacitive loads \cite{Xu2017}. Furthermore, our companion submission \cite{Guerrero2026} has demonstrated repetition rates of 100 kHz, using a poled fibre modulator with smaller V$_\pi$ and shorter pulse widths, thus confirming the potential of the device for practical quantum information processing tasks.

The results obtained are also the first to use poled fibre electro-optical modulation technology for applications in quantum information processing. The fast speed, low insertion loss and direct optical fibre compatibility make it very attractive for demanding applications such as device-independent quantum communication implementations, where low losses and high speeds for application of different measurement settings are crucial \cite{Xavier2025}. Furthermore, the polarisation independence of the device opens further possibilities, such as, in higher-dimensional hybrid entanglement processing where polarisation is one of the degrees-of-freedom \cite{Zhong2024}.



\section*{Acknowledgements}
We acknowledge helpful discussions with Alvaro Alarcón, Walter Margulis, Sebastian Etcheverry, Gustavo do Amaral and Guilherme Tempor\~ao. We acknowledge financing from Vinnova (project no. 2023-01358), the Wallenberg Center for Quantum Technologies (WACQT) and the project “DIGITAL-2022-QCI-02-DEPLOY-NATIONAL” (Project number: 101113375 - NQCIS National Quantum Communication in Sweden) funded by the European Union together with Vinnova and the Wallenberg Centre for Quantum Technology (WACQT).





%\bibliographystyle{unsrt} 

\clearpage
\bibliography{sample.bib}

\end{document}
